\begin{table}[htbp]
\centering
\small
\caption{Teto do Problema D: acurácia da métrica diagonal ajustada ao rótulo real da meia-lua, em R$^2$, R$^3$ e R$^4$.}
\label{tab:bloco1_D_teto_oraculo_meialua}
\begin{tabular}{llrrl}
\toprule
Espaço & Algoritmo & Acurácia vs real & Erros (de 150) & W \\
\midrule
R$^2$ & perceptron & 85,6 (84,7–86,7) & 21,7 & [0,18; 0,95] \\
R$^2$ & nnls & 86,7 & 20,0 & [0,18; 0,79] \\
R$^3$ & perceptron & 86,7 (86,0–87,3) & 20,0 & [0,20; 0,93; 1,01] \\
R$^3$ & nnls & 86,7 & 20,0 & [0,18; 0,79; 0,00] \\
R$^4$ & perceptron & 86,2 (85,3–87,3) & 20,7 & [0,21; 1,03; 0,05; 1,01] \\
R$^4$ & nnls & 86,4 (86,0–86,7) & 20,3 & [0,18; 0,79; 0,00; 0,00] \\
\bottomrule
\multicolumn{5}{p{0.95\linewidth}}{\footnotesize W ajustado diretamente ao rótulo real da meia-lua (150 pontos): teto da métrica diagonal no Problema D. R$^3$ acrescenta x1·x2, R$^4$ acrescenta x1², x2². Médias das sementes (mín–máx).} \\
\end{tabular}
\end{table}
